%================================
% note-setup-leftsidebox.tex
% fenglielie@qq.com 2025-09-12
%================================

\usepackage{amsmath,amsthm,amsfonts,amssymb}
\usepackage{mathtools}
\usepackage{mathrsfs}
\usepackage{bm}
\usepackage{extarrows}
\usepackage[a4paper, margin=1in]{geometry}
\usepackage{float}
\usepackage{indentfirst}
\usepackage{anyfontsize}
\usepackage{booktabs,multirow,multicol}
\usepackage[shortlabels,inline]{enumitem}
\usepackage{appendix}

\usepackage[dvipsnames]{xcolor}
\usepackage{graphicx}
\graphicspath{
    {./figure/}{./figures/}{./image/}{./images/}{./graphic/}{./graphics/}{./picture/}{./pictures/}
}
\usepackage{subcaption}

\usepackage[ruled,linesnumbered,noline]{algorithm2e}
\usepackage{listings}
\lstdefinestyle{simpleStyle}{
    basicstyle=\ttfamily\small,
    breaklines=true,
    keywordstyle=\color{blue},
    identifierstyle=\color{black},
    stringstyle=\color{violet},
    commentstyle=\color[RGB]{34,139,34},
    showstringspaces=false,
    numbers=left,
    numbersep=2em,
    numberstyle=\footnotesize,
    frame=single,
    framesep=1em,
}
\lstset{style=simpleStyle}

\usepackage{imakeidx}  % 用于生成索引 

\usepackage{hyperref}
\hypersetup{
    colorlinks=true,linkcolor=,urlcolor=cyan
}

\renewcommand*{\proofname}{\itshape 证明}  % 证明环境开头改为斜体“证明”

\usepackage{thmtools}

% 注记样式
\declaretheoremstyle[
  headformat={\NAME~{\itshape\NUMBER}\NOTE},  % 编号单独用斜体
  headpunct={},                                  % 标题末尾的标点
  headfont=\itshape,                             % “注记”二字斜体
  bodyfont=\itshape,                              % 正文斜体
  postheadspace=0.5em,                            % 标题与正文间距
]{remarkstyle}

% 重新定义 plain 样式，去除句点
\declaretheoremstyle[
  headpunct={},                         % 去掉编号后的句点
  headformat={\NAME~\NUMBER\NOTE},      % 格式：名称 + 编号（无标点）
  headfont=\bfseries,                   % 标题加粗（plain 默认）
  bodyfont=\itshape,                    % 正文斜体（plain 默认）
  postheadspace=0.7em,                  % 标题后间距
  numbered=yes,                         % 默认编号（但环境仍可覆盖）
]{plain}

% 重新定义 definition 样式，去除句点
\declaretheoremstyle[
  headpunct={},
  headformat={\NAME~\NUMBER\NOTE},
  headfont=\bfseries,
  bodyfont=\normalfont,                 % 正文正体（definition 默认）
  postheadspace=0.7em,
  numbered=yes,
]{definition}

%% define environments

\declaretheorem[style=plain, name=定理, numbered=yes, numberwithin=section]{theorem}
\declaretheorem[style=plain, name=定理, numbered=no]{theorem*}

\declaretheorem[style=plain, name=命题, numbered=yes, sibling=theorem]{proposition}
\declaretheorem[style=plain, name=命题, numbered=no]{proposition*}

\declaretheorem[style=plain, name=推论, numbered=yes, sibling=theorem]{corollary}
\declaretheorem[style=plain, name=推论, numbered=no]{corollary*}

\declaretheorem[style=plain, name=引理, numbered=yes, sibling=theorem]{lemma}
\declaretheorem[style=plain, name=引理, numbered=no]{lemma*}

\declaretheorem[style=plain, name=断言, numbered=yes, sibling=theorem]{claim}
\declaretheorem[style=plain, name=断言, numbered=no]{claim*}

\declaretheorem[style=definition, name=定义, numbered=yes, numberwithin=section]{definition}
\declaretheorem[style=definition, name=定义, numbered=no]{definition*}

\declaretheorem[style=definition, name=例, numbered=yes, numberwithin=section]{example}
\declaretheorem[style=definition, name=例, numbered=no]{example*}

\declaretheorem[style=definition, name=Problem, numbered=yes, numberwithin=section]{problem}
\declaretheorem[style=definition, name=Problem, numbered=no]{problem*}

\declaretheorem[style=remarkstyle, name=注记, numbered=yes, numberwithin=section]{remark}
\declaretheorem[style=remarkstyle, name=注记, numbered=no]{remark*}

\declaretheoremstyle[headfont=\color{orange!80}\bfseries, bodyfont=\normalfont, spaceabove=3pt, spacebelow=3pt]{notestyle}

\declaretheorem[style=notestyle, name=Note, numbered=yes, numberwithin=section]{note}
\declaretheorem[style=notestyle, name=Note, numbered=no]{note*}

\declaretheoremstyle[headfont=\bfseries, bodyfont=\normalfont, spaceabove=3pt, spacebelow=3pt, qed=\ensuremath{\square}]{solutionstyle}

\declaretheorem[style=solutionstyle, name=Solution, numbered=yes, numberwithin=section]{solution}
\declaretheorem[style=solutionstyle, name=Solution, numbered=no]{solution*}

\usepackage[most]{tcolorbox}

\newcommand{\newtcbenvironment}[2]{
    \tcolorboxenvironment{#1}{#2, enhanced, breakable, sharp corners,leftrule=2pt, rightrule=0pt, toprule=0pt, bottomrule=0pt}
    \tcolorboxenvironment{#1*}{#2, enhanced, breakable, rounded corners,leftrule=2pt, rightrule=0pt, toprule=0pt, bottomrule=0pt}
}

%% define styles

\newtcbenvironment{theorem}{colframe=RoyalPurple, colback=RoyalPurple!8}
\newtcbenvironment{proposition}{colframe=RoyalPurple, colback=RoyalPurple!8}
\newtcbenvironment{corollary}{colframe=NavyBlue, colback=SkyBlue!8}
\newtcbenvironment{lemma}{colframe=NavyBlue, colback=SkyBlue!8}
\newtcbenvironment{claim}{colframe=NavyBlue, colback=SkyBlue!8}

\newtcbenvironment{definition}{colframe=ForestGreen, colback=ForestGreen!5}
\newtcbenvironment{example}{colframe=RawSienna, colback=RawSienna!5}
\newtcbenvironment{problem}{colframe=WildStrawberry!30, colback=WildStrawberry!5}

%% cbox
\newtcolorbox{cbox}[1][]{%
    enhanced,
    breakable,
    sharp corners,
    leftrule=2pt, rightrule=0pt, toprule=0pt, bottomrule=0pt,
    colframe=SkyBlue,
    colback=SkyBlue!8,
    #1
}

%% cover
\usepackage{titling}
\newcommand{\extrainfo}{}
\renewcommand{\extrainfo}[1]{\renewcommand{\extrainfocontent}{#1}}
\newcommand{\extrainfocontent}{}
\newcommand{\makecover}[1]{%
    \begin{titlepage}
        \newgeometry{margin=0in}
        \parindent=0pt
        \includegraphics[width=\linewidth]{#1} % size = 1280*1024
        \vfill
        \begin{center}
            \parbox{0.618\textwidth}{%
                \raggedleft{\bfseries \Huge \thetitle} \\[0.6pt]
                \rule{0.618\textwidth}{4pt} \\
            }
        \end{center}
        \vfill
        \begin{center}
            \parbox{0.618\textwidth}{%
                \raggedleft\Large
                \begin{tabular}{r}
                    \theauthor \\
                    \thedate   \\
                \end{tabular}%
            }
        \end{center}
        \vfill
        \begin{center}
            \parbox[t]{0.7\textwidth}{\centering \itshape \extrainfocontent}
        \end{center}
        \vfill
    \end{titlepage}
    \restoregeometry
    \thispagestyle{empty}
}
% USAGE
% \extrainfo{xxx}
% \makecover{/path/to/cover.png}
